\documentclass[11pt]{article}
\usepackage[utf8]{inputenc}
\usepackage{amsmath,amssymb}
\usepackage{geometry}
\geometry{margin=1in}
\usepackage{hyperref}
\usepackage{booktabs}

\title{AURA-100: An Open Reference Design for Post-Aural Wearable\\
Brain-Computer Interfaces with g3-Honest Regulatory Scaffolding}
\author{dancinlab AURA project}
\date{2026}

\begin{document}
\maketitle

\begin{abstract}
The wearable EEG-based brain-computer interface (BCI) ecosystem is
fragmented: open-source tooling excels at signal analysis (MNE-Python,
EEGLAB) yet lags in acquisition hardware and clinical-grade regulatory
scaffolding, where proprietary platforms dominate. We present AURA-100,
an open hexa-native reference design for a post-aural / temporal-bone
clip wearable BCI, organized around a seven-verb pipeline
(specify $\to$ structure $\to$ design $\to$ analyze $\to$ synthesize
$\to$ verify $\to$ handoff). Our contributions are: (1) an open
reference design with illustrative bill-of-materials; (2) honest
regulatory dossier scaffolding spanning FDA 510(k) Class II, EU MDR
Class IIa, IEC 60601-1-2 EMC, and Bluetooth SIG qualification; (3) a
cross-vendor SoC and Lab Streaming Layer (LSL) bridge proposal
(Nordic nRF5340, Espressif ESP32-C6, OpenBCI); and (4) an honest-gap
declaration for MRI-safety, where Sim4Life holds an FDA-MDDT monopoly
and we sketch a multi-year open-source breakthrough path. A key
numerical result: a naive-DFT alpha-band periodogram reproduces the
MNE-Welch (boxcar) estimator at float64 machine precision
($\text{rel\_spread} = 2.23 \times 10^{-16}$) on real PhysioNet
EEGBCI data, generalizing the Sleep-EDF spectral parity
(mean\_rel\_err $= 8.4 \times 10^{-7}$). All artifacts are
template / illustrative (\texttt{absorbed=false} where accredited
closure is pending); the \texttt{g3} honesty discipline keeps every
claim ring-fenced.
\end{abstract}

\section{Introduction}
Wearable BCI landscape; open-vs-proprietary gap; the case for an open
post-aural reference design.

\section{Seven-Verb Design Pipeline}
V0--V7 producer cells and architectural choices. V0 (verify) is the
sole \texttt{absorbed=true} measured-oracle cell, anchored on
PhysioNet Sleep-EDF spectral parity.

\section{Illustrative Substance (D1--D7)}
Bill-of-materials (Texas Instruments ADS1299, Nordic nRF5340, Li-Po),
ngspice AFE noise modeling, openEMS antenna / SAR, MNE-Python band
power. Every value is an illustrative reference target.

\section{Breakthrough Paths (G1--G3)}
G1 open MRI-safety solver (ASTM F2182, ISO 10974) at Sim4Life parity;
G2 open EEG acquisition platform; G3 RISC-V BLE SoC diversity.

\section{Regulatory Scaffolding (A1--A4)}
FDA 510(k) eSTAR, EU MDR Annex IX, IEC 60601-1-2 EMC, Bluetooth SIG
qualification --- submission-ready drafts.

\section{Numerical Verification}
\begin{table}[h]
\centering
\begin{tabular}{lll}
\toprule
Dataset & Estimator pair & Agreement \\
\midrule
Sleep-EDF & naive-DFT vs MNE-Welch & mean\_rel\_err $8.4\times10^{-7}$ \\
EEGBCI s1r1 & naive-DFT vs Welch(Hann) & rel\_err $0.28$ (window mismatch) \\
EEGBCI s1r1 & naive-DFT vs Welch(boxcar) & rel\_spread $2.23\times10^{-16}$ \\
\bottomrule
\end{tabular}
\caption{Alpha-band periodogram parity. Window convention (Hann vs
no-window) is the only variable; with boxcar pinned, the naive-DFT and
MNE-Welch estimators agree to float64 machine precision.}
\end{table}

\section{g3 Honesty Discipline}
Scope-caveat methodology; \texttt{absorbed=false} rationale; the
distinction between in-silico closure and accredited-lab confirmation.

\section{Discussion and Limitations}
Open ecosystem implications; Sim4Life monopoly; RISC-V medical
silicon; template-substance limitations and downstream gates.

\section{Conclusion}
hexa-native open reference design as an honest scaffolding pattern for
regulated wearable BCI development.

\end{document}
